\documentclass[a4paper]{article}

\title{24(?) Preludes and Fugues:\\an anthology of writing}
\author{Sadie}

\newcommand\fl{$\flat$}
\newcommand\sh{$\sharp$}

\begin{document}
\maketitle

This is a planned anthology of good writing with a few unique features.
It is based around the structure of J.~S.~Bach's 48 preludes and
fugues for keyboard but note that \textbf{you do not need to know
  anything about music to be able to contribute.}

If you wish to take part please read at least the section here marked
`Essential' and email Sadie at \texttt{sadie@sadie.org.uk} ASAP with
requested information.  The deadline for final copy is March 31st,
2026.

\section{Essential}

The proposal here is to produce an anthology of a number (ideally 24
but more or less should be possible) of pieces of original writing
based on Bach's `Preludes and Fugues.'  Each piece should be
structured as a short Introduction (the prelude) followed by a Story
(the fugue) with a page-break between.  The entire piece (prelude
and fugue) should feature a particular mood, emotion or trait
analogous to the musical concept of \textit{key}.  (There are 24
keys in classical music.)

The book will be published on Amazon.  Authors will retain copyright
of their own work and can re-publish it later, in a collection
perhaps. Authors will receive one copy of the collection and can
buy additional copies at cost price.

To contribute, this is what you must do.

\begin{itemize}
\item Decide on a small number (from three to six) of words describing
  a mood, emotion or character trait you wish to addess.  One
  suggestion is that you might choose these from listening to some music
  and recording your emotional responses to it 
  but you can come up with the mood in
  any way you like. From now on your work will focus on these words
  and not the music if you used any.
\item Let Sadie know ASAP that you want to take part and tell her
  the words you have chosen.
\item Decide on a place, object, animal or person which represents
  this mood, character, emotion or traits.  This is your Subject.  (It
  is best if the subject is not musical or a piece of music but
  something quite different, though your writing may involve music if
  you wish.)
\item Write a short introduction, introducing your subject and setting
  the scene and/or providing a short back-story (the prelude). The
  prelude sets the scene but is not the main story.
\item Write a longer story about the subject (the fugue). 
\item Submit for consideration ASAP and before the deadline of March 31, 2026.
\end{itemize}

I intend `story' to be taken in the most general possible sense.  It
may help you (but it is not essential) to know that `Fugue' literally
means `Flight' as in `running away.'

Your work may be in any medium: prose, poetry, play, dialogue,
monologue or journalism.  It may be fact or fiction.  I don't want to
specify a length here other than saying that over-long pieces will not
be accepted.  Please focus your writing and not be repetitive or use
more words than necessary.  For prose or drama, 6-10 pages, up to about 3500
words, is appropriate though shorter or much shorter is fine.  (Please
discuss if this does not fit your vision.) For poetry, up to about 200
lines is fine, possibly more.

You should not engage in any inappropriate dismissive statements
against any individuals or minorities.  The anthology is intended to
be inclusive of all kinds of people and your piece will be rejected
if this is not the case.

Submissions should be as plain text documents or Word format
\textbf{without formatting} apart from indications of title(s), bold
font or italics and new paragraphs or new lines and stanzas. (This is
important.) If you have special formatting instructions please detail
them separately in an additional document e.g.~as a separate PDF
document to give me the idea of what you want.  (Do not expect any
formatting in your word document to be correctly transmitted except
for the features just listed.)

\textbf{Please send ASAP an email expressing your wish to take part.}
This must contain,
\begin{itemize}
\item Your name and preferred contact details
\item The name or pseudonym you want your to work to appear under.
\item Your words for the `mood'
\end{itemize}

\textbf{If you change your mind and cannot contribute after all please
  send an email ASAP.}

\textbf{Please send your copy by 31st March 2026 AT THE LATEST.} This should include
\begin{itemize}
\item Name as you'd like it to appear and a short biographical
  statement---up to 100 words.
\item Your words for the `mood'.
\item Your prelude.
\item Your fugue.
\item An acknowledgement that this is all your own work and you hold
  all rights to it and it is not libelous, slanderous or misleading
  etc.
\end{itemize}

For all communications please email Sadie at \texttt{sadie@sadie.org.uk}.

A longer version of this prompt document is available at\[
\texttt{https://sadie.org.uk/preludes-and-fugues}\]

\clearpage

\section{Further prompts, ideas and comments}

This is a suggestion to get started:
\begin{itemize}
\item Listen to ONE piece of music carefully several times.  Focus on
  the mood or emotions it conveys and your response to it. Ideally the
  music should NOT have words or if it does these should be in a
  language you do not understand.
\item Write down 3--6 words describing the mood or emotion or traits you hear.
  Don't make this over-complicated.
\item Think of a Subject: a person, place, animal or object that represents the
  mood or emotion or traits you have just described.  It should NOT be musical
  and in particular it MUST NOT be the music you have just been listening to.
\item Picture your subject and introduce it in its context.  This is your Prelude.
\item Write a story about your subject.  This is your Fugue.
\item Give your piece (that is, the prelude and fugue combined
without the initial adjectives) a title. 
\end{itemize}

Please take the precise instructions here with a pinch of salt if
necessary and email Sadie if you are unclear or want to check if your
great idea can work in this context.  The objective is to have a
common structure to the contributions but within that to have plenty
of variety. I don't want you to write \textit{about} a piece of music
particularly, but to write about a subject which represents the
emotions you feel from the music.

\section{Musical background}

J.~S.~Bach famously wrote 48 preludes and fugues for keyboard.  This
collection is known as the \textit{Well-tempered Klavier} and
colloquially as `The 48.' There are 24 preludes and fugues in book~I,
one in each of the 24 possible \textit{keys} in classical music.
Book~II contains a further 24 preludes and fugues in each of the
\textit{keys}.  (I use italics to denote a technical music term.  It
is not important if you don't know these terms.)

Many people (myself included) believe that each \textit{key} conveys a
\textbf{mood} or \textbf{feeling} unique to itself.  Exactly what that
is, is a matter for debate or argument. No doubt this is a personal
response that might vary with context, style, etc.  Some
\textit{keys} are familiar and used frequently and some are unusual,
and this contributes to the corresponding \textbf{mood}.

If you are not sure what music to listen to, I would suggest listening
to any ONE of the Bach 48 Preludes and Fugues or any ONE of the Chopin
\textit{Preludes}. There is one of these in each key.

If you do do this work by starting with a piece of music, it would be
very interesting to know (a)~which piece inspired you and/or (b)~the key it is in.

\section{Keys}

I list here all the 24 classical \textit{keys}.  Each key has a
\textit{tonic note} (there are 12 of these) and is either
\textit{major} or \textit{minor}.  Each key has an associated number
of \textit{sharps} or \textit{flats} which indicates how it is
written.  Some keys can be written equivalently in either
\textit{sharps} or \textit{flats}.  The more \textit{sharps} or
\textit{flats} a key has the more unusual and esoteric it sounds. Keys
with few sharps or flats sound more `normal.'  I also write a few
words about each key that I have seen written on the internet.  (I
don't agree with them all.)

\medskip

\noindent\textbf{C major.}  No sharps or flats.
Pure, innocence, simplicity, naïvety, children's talk.  Often the first (easiest) key
we learn to play in.

\noindent\textbf{A minor.}  No sharps or flats.
Pious, feminine, tenderness.

\noindent\textbf{G major.}  One sharp.
Rustic, idyllic, lyrical, calm, gentle, peaceful.

\noindent\textbf{E minor.}  One sharp.
Na\"ive, innocent, lament.

\noindent\textbf{D major.}  Two sharps.
Triumphant. Confident, victorious, rejoicing, bright.

\noindent\textbf{B minor.}  Two sharps.
Quiet acceptance of fate. Submission.

\noindent\textbf{A major.}  Three sharps.
Innocent love, satisfaction,
cheerfulness and trust.

\noindent\textbf{F\sh\ minor.}  Three sharps.
Gloomy.
Resentment and discontent.

\noindent\textbf{E major.}  Four sharps.
Fresh, brash, new, happy, joy, laughing, pleasure.

\noindent\textbf{C\sh\ minor.}  Four sharps.
Penitential, lamentation, intimate, sighing.

\noindent\textbf{B major.}  Five sharps.
Colourful, passionate. Can incorporate anger, rage, jealousy, fury, despair.

\noindent\textbf{G\sh\ minor aka A\fl\ minor.}  Five sharps or seven flats.
Grumbling, suffocating, wailing, struggle, difficulty.

\noindent\textbf{F\sh\  major aka G\fl\ major.}  Six sharps or six flats.
Triumph over difficulty.

\noindent\textbf{D\sh\ minor aka E\fl\ minor.}  Six sharps or six flats.
Anxiety, distress,  despair,  depression, gloom.

\noindent\textbf{D\fl\ major aka C\sh\ major.}  Five flats or seven sharps.
Majestic.

\noindent\textbf{B\fl\ minor aka A\sh\ minor.}  Five flats or seven sharps.
Quaint, dressed in a garment of night. Surly and mocking.

\noindent\textbf{A\fl\ major.}  Four flats.
Pomp.  Grand. Possibly: death, grave, judgement, eternity.

\noindent\textbf{F minor.}  Four flats.
Melancholy. Sombre, emotional.

\noindent\textbf{E\fl\ major.}  Three flats.
Considered by many composers as `heroic.' Also devotion.

\noindent\textbf{C minor.}  Three flats.
Tragic, fate, enigmatic, brooding,
emotionally charged, melancholy, solemnity.

\noindent\textbf{B\fl\ major.}  Two flats.
Noble. Cheerful, clear conscience and hope.

\noindent\textbf{G minor.}  Two flats.
Sweetness, tenderness, serious, magnificent, discontent, unease, intensity.

\noindent\textbf{F major.}  One flat. 
Pastoral.  Countryside. Calm.

\noindent\textbf{D minor.}  One flat.
Melancholy, feminine, brooding humours.

\section{Other pieces you could listen to}

Here, I list a few other pieces you can check out that I think
are characteristic of their keys. (There are many others
of course.)  For multi-movement works e.g.~symphonies and concertos,
unless it says otherwise, \textbf{the first and last movements are in the key
specified} and the other movements may not be.

\medskip

\noindent\textbf{C major.}

J.~S.~Bach, prelude in C major. \textit{Well-tempered klavier} book 1 number 1.  This is
very well known indeed.

Beethoven, Symphony number 1.

Shostakovitch,  end of symphony number 7 (Leningrad).

Duke Ellington, \textit{C jam blues}.

Schubert, Great C major symphony.

\noindent\textbf{A minor.}

Grieg, Piano concerto

Mendelssohn symphony number 3 (Scottish)

\noindent\textbf{G major.}

Dvorak, Symphony number 8 (this alternates G major/minor throughout the first and last movements finishing in G major)

Mahler Symphony 4, first movement.

\noindent\textbf{E minor.}

J.~S.~Bach, St Matthew Passion.

Dvorak, Symphony 9. (New World)

Mendelssohn, Violin Concerto.

Brahms Symphony number 4.

\noindent\textbf{D major.}

The Hallelujah chorus from Handel's Messiah.

The last movement to Mahler's first symphony.

Mozart, Prague symphony.

Tchaikovsky, Violin concerto.

\noindent\textbf{B minor.}

J.~S.~Bach, h-Moll Messe.

Dvorak, Cello concerto.

Grieg, In the hall of the mountain King.

\noindent\textbf{A major.}

Beethoven Symphony number 7.

Mozart Clarinet concerto.

Shostakovitch Festive overture

Schubert Trout quintet

\noindent\textbf{F\sh\ minor.}

Haydn Farewell symphony

\noindent\textbf{E major.}

Vivaldi, Four seasons, Concerto number 1, Spring.

Rossini, William Tell overture.

Wagner, Siegfried Idyll.

Grieg, Morning mood.

\noindent\textbf{C\sh\ minor.}

First movement of Beethoven Moonlight sonata

First movement of Mahler 5th symphony

\noindent\textbf{B major.}

Beethoven symphony number 4.

\noindent\textbf{G\sh\ minor aka A\fl\ minor.}

The pieces in this key from the Bach 48 or Chopin preludes

\noindent\textbf{F\sh\  major aka G\fl\ major.}

The pieces in this key from the Bach 48 or Chopin preludes

\noindent\textbf{D\sh\ minor aka E\fl\ minor.}

Prokofiev 6th symphony

Dave Brubeck, `Take five' (when played properly!)

\noindent\textbf{D\fl\ major aka C\sh\ major.}

Debussy, Clair de lune, number 3 of Suite bergamasque

Dvorak Symphony 9 (New world) 2nd movement.

Tchaikovsky Piano concerto no 1, second movement.

\noindent\textbf{B\fl\ minor aka A\sh\ minor.}

Barber, Adagio for strings

Tchaikovsky Piano concerto number 1

Walton Symphony number 1

\noindent\textbf{A\fl\ major.}

Elgar symphony number 1.

\noindent\textbf{F minor.}

Beethoven Egmont overture.  (The end is in a contrasting F major.)

Dukas, the sorcerer's apprentice.

\noindent\textbf{E\fl\ major.}

Beethoven's 3rd symphony, \textit{Eroica}.

Richard Strauss's tone poem \textit{Ein Heldenleben}.

Mozart symphony 39

Tchaikovsky 1812 overture

\noindent\textbf{C minor.}

The first movement of Beethoven's fifth symphony (d-d-d-DUUURH).
(The last movement is in a contrasting C major.)

Brahms, symphony number 1.

Beethoven pathetique sonata, first movement.

\noindent\textbf{B\fl\ major.}

Hindemith Symphony in B-flat for Band 

\noindent\textbf{G minor.}

Mozart symphony number 40

Allegri, Miserere mei

\noindent\textbf{F major.}

Beethoven symphony 6

Beethoven symphony 8

Brahms symphony 3

\noindent\textbf{D minor.}

Franck, Symphony in D minor.

Schumann, Symphony number 4.

Shostakovitch symphony number 5, first movement.

Beethoven Symphony 9, first movement.

\section{Modulation}

Ask if you have other ideas or questions. One advanced suggestion
is to introduce a modulation by introducing a second
subject in a different key. \textbf{But you should return to the
  original subject/key at the end.}  In music, there are some keys that
it is `natural' to modulate to
but I see no particular reason to restrict yourself.

Good luck! \hfill Sadie 

\end{document}
